For $pq\ne6$, this counts interior lattice points of the parallelogram spanned by $(p,2)$ and $(3,q)$. Its area is $|pq-6|$ and its boundary lattice count is $2\gcd(p,2)+2\gcd(3,q)$. Pick\textquotesingle s theorem gives $14=|pq-6|-\gcd(p,2)-\gcd(3,q)+1$. Thus $pq$ lies between 21 and 24. Checking factor pairs gives $(2,11),(2,12),(3,7),(4,6),(8,3),(21,1),(22,1)$, hence $\boxed7$. If $pq=6$, the map is singular and each attainable point has a continuum of preimages, so the count cannot be exactly 14.
\begin{asy}
import graph;
usepackage("amsmath");
import geometry;
size(7cm);
pair O=(0,0),U=(4,2),V=(3,6);draw(O--U--U+V--V--cycle);label("$(0,0)$",O,(-1,-1));label("$(p,2)$",U,(1,-1));label("$(3,q)$",V,(-1,1));label("$(p+3,q+2)$",U+V,(1,1));
\end{asy}

